
This work asked whether auxiliary training objectives could teach language models to preserve user agency. The answer, at proof-of-concept scale, is: not automatically.

BiDPO demonstrates that agency-like signals exist orthogonally in preference data (r = $-0.026)$ and integrate stably into DPO training without numerical instabilities. However, the gap between training-time signals and generation-time behavior remains (p = 0.247, d = 0.016).

Three directions for future work:

\begin{enumerate}
\item \textbf{Full-scale training:} Extend training to a full epoch (~10,000 steps) to determine whether longer optimization induces behavioral change.
\end{enumerate}

\begin{enumerate}
\item \textbf{$\lambda$ sweep:} Test $\lambda \in$ {0.25, 0.5, 0.75, 1.0} to identify whether stronger or weaker agency weighting improves transfer.
\end{enumerate}

\begin{enumerate}
\item \textbf{Learned collaboration scores:} Replace heuristic pattern matching with a learned classifier trained on human annotations of agency-preserving responses.
\end{enumerate}

The gap identified here is likely relevant beyond BiDPO---any auxiliary objective for alignment faces the challenge of transferring training-time gradient pressure to inference-time behavior.

